\documentclass[12pt,a4paper]{report}

% ============================================================
% PACKAGES
% ============================================================
\usepackage[a4paper, margin=2.5cm, top=3cm, bottom=3cm]{geometry}
\usepackage[T1]{fontenc}
\usepackage{lmodern}
\usepackage{microtype}
\usepackage{amsmath}
\usepackage{amssymb}
\usepackage{mathtools}
\usepackage{booktabs}
\usepackage{longtable}
\usepackage{tabularx}
\usepackage{multirow}
\usepackage{array}
\usepackage{graphicx}
\usepackage{xcolor}
\usepackage{listings}
\usepackage{hyperref}
\usepackage{enumitem}
\usepackage{float}
\usepackage{caption}
\usepackage{subcaption}
\usepackage{tcolorbox}
\usepackage{titlesec}
\usepackage{setspace}
\usepackage{mdframed}
\usepackage{fancyhdr}
\usepackage{appendix}
\usepackage{verbatim}

% ---- Colour palette ----
\definecolor{teal}{RGB}{0, 180, 145}
\definecolor{darkslate}{RGB}{10, 20, 35}
\definecolor{codeblue}{RGB}{30, 110, 200}
\definecolor{codegray}{RGB}{246, 248, 252}
\definecolor{codered}{RGB}{190, 45, 45}
\definecolor{amber}{RGB}{240, 158, 35}
\definecolor{bordercolor}{RGB}{200, 215, 230}
\definecolor{accentblue}{RGB}{30, 80, 160}

% ---- Hyperref setup ----
\hypersetup{
  colorlinks=true,
  linkcolor=teal,
  urlcolor=accentblue,
  citecolor=teal,
  pdftitle={Industrial Energy Efficiency Copilot: Full Technical Report},
  pdfauthor={CarbonTatva Engineering},
  pdfsubject={Cloud-Enabled Agentic RAG System — Architecture, Evaluation, Fine-Tuning},
}

% ---- Code listings ----
\lstset{
  backgroundcolor=\color{codegray},
  basicstyle=\footnotesize\ttfamily,
  keywordstyle=\color{codeblue}\bfseries,
  stringstyle=\color{codered},
  commentstyle=\color{gray}\itshape,
  numberstyle=\tiny\color{gray},
  numbers=left,
  stepnumber=5,
  numbersep=8pt,
  frame=single,
  rulecolor=\color{bordercolor},
  breaklines=true,
  breakatwhitespace=true,
  captionpos=b,
  keepspaces=true,
  showspaces=false,
  showstringspaces=false,
  tabsize=4,
  xleftmargin=14pt,
}

% ---- Chapter / section formatting ----
\titleformat{\chapter}[display]
  {\Large\bfseries\color{darkslate}}
  {\color{teal}\chaptertitlename\ \thechapter}{12pt}
  {\huge\bfseries\color{darkslate}}
  [\vspace{2pt}\color{teal}\rule{\textwidth}{1.2pt}]

\titleformat{\section}
  {\large\bfseries\color{darkslate}}{\thesection}{1em}{}
  [\color{bordercolor}\rule{\textwidth}{0.5pt}]

\titleformat{\subsection}
  {\normalsize\bfseries\color{darkslate}}{\thesubsection}{1em}{}

\titleformat{\subsubsection}
  {\normalsize\bfseries}{\thesubsubsection}{1em}{}

% ---- Header / Footer ----
\pagestyle{fancy}
\fancyhf{}
\fancyhead[L]{\small\color{gray}Industrial Energy Efficiency Copilot}
\fancyhead[R]{\small\color{gray}\leftmark}
\fancyfoot[C]{\small\thepage}
\renewcommand{\headrulewidth}{0.4pt}
\setlength{\headheight}{15pt}

% ---- Coloured boxes ----
\tcbuselibrary{skins,breakable}
\newtcolorbox{notebox}[1][Note]{
  enhanced, breakable,
  colback=teal!7!white,
  colframe=teal,
  boxrule=0.5pt,
  arc=4pt,
  title={#1},
  fonttitle=\bfseries\small,
  attach boxed title to top left={yshift=-2pt, xshift=8pt},
}
\newtcolorbox{warnbox}[1][Warning]{
  enhanced,
  colback=amber!10!white,
  colframe=amber!70!black,
  boxrule=0.5pt,
  arc=4pt,
  title={#1},
  fonttitle=\bfseries\small,
  attach boxed title to top left={yshift=-2pt, xshift=8pt},
}
\newtcolorbox{keybox}[1][Key Principle]{
  enhanced,
  colback=accentblue!6!white,
  colframe=accentblue,
  boxrule=0.5pt,
  arc=4pt,
  title={#1},
  fonttitle=\bfseries\small,
  attach boxed title to top left={yshift=-2pt, xshift=8pt},
}

% ---- Spacing ----
\setlength{\parskip}{0.55em}
\setlength{\parindent}{0pt}
\onehalfspacing

% ============================================================
\begin{document}

% ============================================================
% TITLE PAGE
% ============================================================

\begin{titlepage}
\centering
\vspace*{1.8cm}

{\Huge\bfseries Industrial Energy Efficiency Copilot\par}
\vspace{0.6cm}
{\LARGE Full Technical Report\par}
\vspace{0.3cm}
{\large\color{gray}Architecture · Implementation · Evaluation · Future Scope\par}

\vspace{1.2cm}
\noindent\rule{0.85\textwidth}{1.2pt}
\vspace{0.5cm}

{\large
\textbf{System:} LangGraph · Mistral API · Tesseract OCR · Hybrid RAG \\[0.3em]
\textbf{Version:} 3.0 (Production Release) \\[0.3em]
\textbf{Organisation:} CarbonTatva Engineering \\[0.3em]
\textbf{Date:} April 2026 \\[0.3em]
\textbf{Status:} \textcolor{teal}{Implementation Complete · Awaiting Full Ingestion}
}

\vspace{1.2cm}
\noindent\rule{0.85\textwidth}{0.4pt}

\vspace{0.8cm}

\begin{abstract}
\noindent
This document is the complete technical reference for the Industrial Energy Efficiency Copilot---a production-grade, agentic Retrieval-Augmented Generation (RAG) system grounded exclusively in the Bureau of Energy Efficiency (BEE) Thermal and Electrical Utilities manuals. The report covers the full engineering stack: Tesseract-based OCR ingestion, LangGraph orchestration with parallel retrieval, Mistral API cloud inference, Server-Sent Events (SSE) streaming, a streaming Next.js frontend, and a comprehensive automated evaluation framework. The system is designed for high performance with cloud-based inference and public deployment capability. The document also presents a structured fine-tuning strategy as future scope, positioned strictly as a last-resort optimisation after retrieval and prompt improvements have been exhausted.
\end{abstract}

\vfill
{\small\itshape Internal engineering document --- CarbonTatva Engineering. BEE manuals \textcopyright\ Bureau of Energy Efficiency, Government of India.}
\end{titlepage}

\tableofcontents
\listoftables
\newpage

% ============================================================
\chapter{System Overview}
% ============================================================

\section{Motivation}

The Bureau of Energy Efficiency (BEE) publishes two comprehensive manuals for industrial energy professionals:

\begin{itemize}
  \item \textit{Energy Efficiency in Thermal Utilities} --- covers boilers, furnaces, steam systems, combustion, cogeneration, and waste heat recovery.
  \item \textit{Energy Efficiency in Electrical Utilities} --- covers motors, variable frequency drives, power factor correction, tariffs, lighting, and compressed air.
\end{itemize}

Both manuals are fully image-scanned (zero machine-readable text), densely technical, and structured as reference documents rather than browsable resources. Industrial engineers and energy auditors typically spend significant time manually searching these manuals for specific guidance.

The copilot automates this search and synthesises grounded, professional answers using eight specialist tools, leveraging the Mistral API for highly responsive and accurate generation.

\section{Design Goals}

\begin{table}[H]
\centering
\caption{System Design Goals}
\begin{tabular}{@{}lll@{}}
\toprule
\textbf{Goal} & \textbf{Design Choice} & \textbf{Result} \\
\midrule
Cloud Scalability & Mistral API inference & High availability + low TTFT \\
Image PDF support & Tesseract OCR with caching & Full text extraction \\
Structured reasoning & LangGraph pipeline & Deterministic routing \\
Low latency & Parallel retrieval + streaming & TTFT $\approx$ 1--3 s \\
Grounded answers & RAG-only (no hallucination mode) & Clean, professional responses \\
Public Access & Configured for cloud deployment & CORS enabled, scalable \\
\bottomrule
\end{tabular}
\end{table}

\section{System Architecture}

\begin{figure}[H]
\centering
\fbox{%
\begin{minipage}{0.94\textwidth}
\ttfamily\small
\textbf{Frontend} (Next.js 14, SSE streaming) \\
\quad $\downarrow$ POST /api/stream \\
\textbf{FastAPI Backend} \\
\quad $\downarrow$ \\
\textbf{LangGraph Pipeline} (10 nodes) \\
\quad normalize\_query \\
\quad $\downarrow$ \\
\quad planner\_router [Mistral format=json] \\
\quad $\downarrow$ \quad [fan-out] \\
\quad retrieve\_dense [ChromaDB HNSW] \quad \| \quad retrieve\_sparse [BM25] \\
\quad $\downarrow$ \quad [fan-in] \\
\quad merge\_results [RRF fusion] \\
\quad $\downarrow$ \\
\quad rerank [cross-encoder/ms-marco] \\
\quad $\downarrow$ \\
\quad tool\_dispatch [prompt selector] \\
\quad $\downarrow$ \\
\quad answer\_generation [Mistral API streaming, mistral-small-latest] \\
\quad $\downarrow$ \\
\quad evidence\_assembly $\rightarrow$ response\_formatter \\
\textbf{Knowledge Base} (ChromaDB + BM25, local disk) \\
\quad Built from: BEE Thermal + Electrical PDFs via Tesseract OCR
\end{minipage}%
}
\caption{End-to-end system data flow}
\end{figure}

\section{Tool Modes}

\begin{table}[H]
\centering
\caption{Eight Specialist Tool Modes}
\begin{tabularx}{\textwidth}{@{}llX@{}}
\toprule
\textbf{Tool} & \textbf{Query Type} & \textbf{Output Style} \\
\midrule
\texttt{qa} & Direct factual questions & Concise, professional paragraph \\
\texttt{explainer} & Conceptual explanations & Structured sections, progressive depth \\
\texttt{troubleshoot} & Fault diagnosis & Root causes + corrective actions \\
\texttt{opportunity} & Energy saving discovery & Ranked bullet list with estimates \\
\texttt{comparison} & Side-by-side analysis & Tabular comparison with tradeoffs \\
\texttt{checklist} & Audit procedures & Numbered inspection steps \\
\texttt{navigation} & Finding content & Chapter/section/page references \\
\texttt{summarize} & Topic summaries & High-level structured overview \\
\bottomrule
\end{tabularx}
\end{table}

% ============================================================
\chapter{OCR Ingestion Pipeline}
% ============================================================

\section{Problem Statement}

Both BEE PDFs are fully image-scanned with \textbf{zero machine-readable characters} from native PDF text extraction (verified with PyMuPDF). Any system relying on native text extraction indexes empty content and produces zero meaningful answers. Tesseract OCR is therefore mandatory, not optional.

\section{OCR Engine (\texttt{ocr\_engine.py})}

\subsection{Rendering}

Each PDF page is rendered to a PIL Image using PyMuPDF at a configurable DPI (default: 200). Higher DPI improves OCR accuracy at the cost of memory and processing time.

\begin{equation}
  \text{Image width} = \text{page width (pts)} \times \frac{\text{DPI}}{72}
\end{equation}

At 200 DPI the typical page renders to approximately $1650 \times 2200$ pixels.

\subsection{Preprocessing}

Before passing to Tesseract, images are converted to grayscale and mildly sharpened using a convolutional kernel:
\begin{equation}
  K = \begin{pmatrix} 0 & -1 & 0 \\ -1 & 5 & -1 \\ 0 & -1 & 0 \end{pmatrix}
\end{equation}

This improves character edge contrast without introducing artefacts.

\subsection{Dual-PSM Strategy}

Tesseract page segmentation mode (PSM) is applied in two passes:
\begin{enumerate}
  \item \texttt{--psm 1}: Automatic page segmentation with OSD (orientation and script detection).
  \item \texttt{--psm 3} (fallback): Fully automatic page segmentation, triggered when PSM 1 confidence falls below threshold.
\end{enumerate}

The higher-confidence result is retained. Confidence is computed as the mean of per-word confidence scores from \texttt{image\_to\_data()}.

\subsection{OCR Result Object}

\begin{lstlisting}[language=Python, caption={OCR result dataclass}]
@dataclass
class OCRResult:
    page_num:     int
    text:         str      # Extracted text
    confidence:   float    # Mean word confidence (0-100)
    word_count:   int
    ocr_time_ms:  float
    has_content:  bool     # confidence > 40 and word_count > 10
\end{lstlisting}

\section{OCR Cache (\texttt{ocr\_cache.py})}

OCR is computationally expensive: $\sim$4--8 seconds per page on CPU. For a 636-page PDF this is $\sim$60--90 minutes. The cache ensures this cost is paid only once.

\subsection{Cache Key Design}

The cache key is a combination of document ID and PDF file fingerprint:
\begin{equation}
  \text{cache\_key} = \text{doc\_id} + \texttt{\_} + \text{SHA-256}(\text{PDF file})[0:16]
\end{equation}

SHA-256 of the PDF bytes detects any file modification. If the PDF changes, the cache is automatically invalidated.

\subsection{Cache Format}

Each cache file is a JSONL file (one JSON object per page):

\begin{lstlisting}[caption={OCR cache entry}]
{"page":0,"text":"...","confidence":87.4,"word_count":312,"ocr_time_ms":4200}
{"page":1,"text":"...","confidence":91.2,"word_count":298,"ocr_time_ms":3800}
\end{lstlisting}

\section{PDF OCR Loader (\texttt{pdf\_ocr\_loader.py})}

The \texttt{PDFOCRLoader} is a drop-in replacement for the v1 \texttt{PDFLoader}:

\begin{enumerate}
  \item Open PDF with PyMuPDF (\texttt{fitz})
  \item Check cache for all pages --- load from cache if available
  \item For pages not cached: render $\rightarrow$ preprocess $\rightarrow$ OCR $\rightarrow$ cache
  \item Map OCR results to the existing \texttt{ParsedDocument} schema
  \item Return \texttt{ParsedDocument} with \texttt{.pages} containing OCR text
\end{enumerate}

The downstream structure parser and chunker are unchanged --- they operate on \texttt{ParsedDocument} regardless of how the text was extracted.

\section{Ingestion Pipeline (\texttt{scripts/ingest.py})}

\begin{table}[H]
\centering
\caption{Ingestion Pipeline Stages}
\begin{tabular}{@{}rlll@{}}
\toprule
\textbf{Step} & \textbf{Stage} & \textbf{Time (est.)} & \textbf{Output} \\
\midrule
1 & OCR (Tesseract, 200 DPI) & 60--90 min & \texttt{data/ocr\_cache/*.jsonl} \\
2 & Structure parse & $<$1 min & chapter/section hierarchy \\
3 & Multi-granular chunking & $<$2 min & list of \texttt{DocumentChunk} \\
4 & Dense embedding (MiniLM) & 5--15 min & vectors \\
5 & ChromaDB indexing & $<$2 min & \texttt{data/indexes/chroma/} \\
6 & BM25 build & $<$1 min & \texttt{data/indexes/bm25\_index.pkl} \\
\midrule
  & \textbf{Total (first run)} & \textbf{75--110 min} & \\
  & \textbf{Total (subsequent)} & \textbf{10--20 min} & (OCR reused from cache) \\
\bottomrule
\end{tabular}
\end{table}

\section{Chunking Strategy}

Four chunk types are produced from each document:

\begin{itemize}
  \item \textbf{Section chunks} --- full text of a section, preserving headers.
  \item \textbf{Semantic chunks} --- fixed-size windows (512 tokens, 64-token stride) within sections for dense retrieval.
  \item \textbf{Table chunks} --- detected table blocks with structural metadata.
  \item \textbf{List chunks} --- enumerated and bulleted lists extracted as standalone units.
\end{itemize}

Each chunk carries rich metadata: document ID, book name, domain, chapter title, section title, page range, chunk type, and ontology tags.

% ============================================================
\chapter{LangGraph Orchestration Pipeline}
% ============================================================

\section{Motivation for LangGraph}

\begin{table}[H]
\centering
\caption{LangGraph vs Custom Router}
\begin{tabularx}{\textwidth}{@{}lXX@{}}
\toprule
\textbf{Concern} & \textbf{Custom Router (v1)} & \textbf{LangGraph (v2)} \\
\midrule
State management & Implicit via function arguments & Typed \texttt{CopilotState} TypedDict \\
Parallel retrieval & Sequential & Native fan-out/fan-in \\
Per-node latency & One aggregate timer & Per-node timestamps \\
Debug tracing & Print statements & \texttt{debug\_trace} field \\
Error handling & Scattered try/catch & Per-node with graph fallback \\
Testability & Tight coupling & Stateless nodes, mock injection \\
\bottomrule
\end{tabularx}
\end{table}

\section{CopilotState}

All nodes read from and write to a shared typed state:

\begin{lstlisting}[language=Python, caption={CopilotState --- selected fields}]
class CopilotState(TypedDict, total=False):
    user_query:          str
    normalized_query:    str
    session_id:          str
    explanation_level:   str
    domain_hint:         str

    planner:             dict   # Structured planner JSON output
    planner_raw:         str    # Raw JSON string (for debug panel)
    planner_error:       bool   # True if rule-based fallback was used

    dense_candidates:    list[dict]
    sparse_candidates:   list[dict]
    merged_candidates:   list[dict]
    reranked_chunks:     list[dict]

    system_prompt:       str
    answer_draft:        str
    citations:           list[dict]
    final_response:      dict

    debug_trace:         list[str]    # ["normalize_query", "planner_router", ...]
    latency:             dict[str, float]  # node_name -> ms
\end{lstlisting}

\section{Graph Topology}

The graph has 10 nodes connected in a directed acyclic structure with one parallel fan-out:

\begin{enumerate}
  \item \textbf{normalize\_query} --- lowercase, strip, normalise punctuation.
  \item \textbf{planner\_router} --- Mistral API JSON-mode LLM call; rule-based fallback on failure.
  \item \textbf{retrieve\_dense} --- ChromaDB HNSW search, top-20 candidates. \textit{(parallel)}
  \item \textbf{retrieve\_sparse} --- BM25 search, top-20 candidates. \textit{(parallel)}
  \item \textbf{merge\_results} --- Reciprocal Rank Fusion of dense+sparse.
  \item \textbf{rerank} --- Cross-encoder scoring, top-8 retained.
  \item \textbf{tool\_dispatch} --- Selects system prompt from planner decision.
  \item \textbf{answer\_generation} --- Mistral API streaming LLM call.
  \item \textbf{citation\_assembly} --- Maps chunks to citation schema.
  \item \textbf{response\_formatter} --- Builds final response dict.
\end{enumerate}

Nodes 3 and 4 run concurrently via LangGraph's native fan-out, reducing retrieval wall time to $\max$(dense latency, sparse latency) rather than their sum.

\section{Planner Design}

\subsection{Structured JSON Output}

The planner sends a compact prompt to the Mistral API with JSON mode enabled, guaranteeing valid JSON output:

\begin{lstlisting}[caption={Planner JSON output schema}]
{
  "query_type":       "troubleshoot",
  "utility_domain":   "thermal",
  "equipment_tags":   ["boiler", "steam_trap"],
  "candidate_tools":  ["troubleshoot", "qa"],
  "selected_tool":    "troubleshoot",
  "confidence":       0.88,
  "why_selected":     "Diagnostic language with performance drop symptom",
  "retrieval_profile":"semantic",
  "response_style":   "structured_diagnosis"
}
\end{lstlisting}

\subsection{Rule-Based Fallback}

If the LLM planner fails (timeout, parse error, Mistral API unavailable), a deterministic rule-based classifier runs in $<1$ ms. It matches keyword patterns from the query to tool and domain labels with no LLM dependency.

\begin{table}[H]
\centering
\caption{Rule-Based Planner Keyword Patterns}
\begin{tabular}{@{}ll@{}}
\toprule
\textbf{Tool} & \textbf{Trigger patterns} \\
\midrule
\texttt{troubleshoot} & why, drops, failing, problem, issue, root cause \\
\texttt{checklist} & checklist, audit, inspect, verify, commissioning \\
\texttt{comparison} & compare, versus, vs, difference, better, worse \\
\texttt{opportunity} & saving, reduce, improve, conserve, opportunity \\
\texttt{explainer} & explain, how does, what is, describe, mechanism \\
\texttt{summarize} & summarise, summarize, overview, key points \\
\texttt{navigation} & where, which chapter, find, locate, section \\
\texttt{qa} & (default for direct factual questions) \\
\bottomrule
\end{tabular}
\end{table}

\section{Hybrid Retrieval}

\subsection{Dense Retrieval}

Query embedding is produced by \texttt{all-MiniLM-L6-v2} (384-dimensional, runs on CPU, $\sim$15 ms). ChromaDB performs HNSW ANN search against $\sim$10,000 indexed chunks, returning top-20 candidates with metadata filters applied for domain if specified.

\subsection{BM25 Sparse Retrieval}

The query is tokenised and scored against the BM25 index (rank-bm25 library, Okapi BM25 with $k_1 = 1.5$, $b = 0.75$). BM25 excels at exact keyword matches and acronym lookup where dense retrieval may fail.

\subsection{Reciprocal Rank Fusion}

Dense and sparse candidate lists are merged with RRF:

\begin{equation}
  \text{RRF}(d) = \sum_{i \in \{\text{dense}, \text{sparse}\}} \frac{1}{k + \text{rank}_i(d)}, \quad k = 60
\end{equation}

The constant $k=60$ is the standard RRF hyperparameter limiting the influence of very high-ranked documents. No score normalisation is required.

\subsection{Cross-Encoder Reranking}

The merged top-30 candidates are passed to a cross-encoder (\texttt{cross-encoder/ms-marco-MiniLM-L-6-v2}). Unlike the bi-encoder, the cross-encoder attends jointly to query and passage, producing superior relevance scores. The top-6 to 8 reranked chunks are retained for generation.

\section{Mistral LLM Client (\texttt{llm\_client.py})}

The \texttt{MistralClient} exposes two primary modes for cloud inference:

\begin{description}
  \item[\texttt{generate()}] Async full response (for planner JSON mode).
  \item[\texttt{stream()}] Async token generator (for answer streaming).
\end{description}

Streaming is implemented by parsing Server-Sent Events (SSE) from the API:

\begin{lstlisting}[language=Python, caption={Streaming token extraction}]
async for token in self._client.stream(system_prompt, user_prompt, max_tokens, temperature):
    yield token
\end{lstlisting}

% ============================================================
\chapter{FastAPI Backend and SSE Streaming}
% ============================================================

\section{API Endpoints}

\begin{table}[H]
\centering
\caption{Backend API Endpoints}
\begin{tabular}{@{}llll@{}}
\toprule
\textbf{Method} & \textbf{Path} & \textbf{Description} & \textbf{Version} \\
\midrule
GET & \texttt{/api/health} & Health + index status & v1 \\
GET & \texttt{/api/status} & Detailed config dump & v1 \\
POST & \texttt{/api/query} & Full pipeline, single response & v1 \\
POST & \texttt{/api/stream} & Full pipeline, SSE streaming & v2 \\
POST & \texttt{/api/classify} & Routing only, no generation & v1 \\
\bottomrule
\end{tabular}
\end{table}

\section{SSE Streaming Protocol}

The \texttt{/api/stream} endpoint (\texttt{api/stream.py}) emits five event types:

\begin{lstlisting}[caption={SSE event sequence}]
data: {"type":"status","content":"Analysing query..."}

data: {"type":"planner","planner":{
  "selected_tool":"troubleshoot",
  "utility_domain":"thermal",
  "confidence":0.88,...
},"error":false}

data: {"type":"token","content":"Based on the retrieved"}
data: {"type":"token","content":" sections, the root causes are:"}
... (N token events)

data: {"type":"done","citations":[...],"latency_ms":7800,
  "tool_mode":"troubleshoot","retrieval_count":6,
  "follow_up_suggestions":["Check excess air...","..."]}

data: [DONE]
\end{lstlisting}

The planner event arrives before any tokens, allowing the frontend to immediately display the routing decision in the debug panel while generation is still in progress.

\section{Query Cache}

The \texttt{QueryCache} (\texttt{cache/query\_cache.py}) uses \texttt{diskcache} for a persistent, LRU-evicting disk-backed cache:

\begin{itemize}
  \item \textbf{Key:} SHA-256(query + tool\_mode + domain + explanation\_level)
  \item \textbf{TTL:} 1 hour
  \item \textbf{Hit behaviour:} Return cached JSONL events directly (sub-millisecond)
  \item \textbf{Location:} \texttt{data/query\_cache/}
\end{itemize}

% ============================================================
\chapter{Frontend}
% ============================================================

\section{Technology}

The frontend is a Next.js 14 App Router application with:
\begin{itemize}
  \item Vanilla CSS dark design system (no Tailwind dependency)
  \item \texttt{react-markdown} + \texttt{remark-gfm} for structured answer rendering
  \item Custom SSE fetch reader in TypeScript (\texttt{lib/api.ts})
  \item Zero state management library (React hooks only)
\end{itemize}

\section{Streaming Display}

Token events from \texttt{/api/stream} are appended to the message state directly:

\begin{lstlisting}[language=JavaScript, caption={SSE token accumulation}]
let fullText = '';
onToken: (token) => {
  fullText += token;
  setMessages(prev => prev.map(m =>
    m.id === assistantId
      ? { ...m, content: fullText }
      : m
  ));
}
\end{lstlisting}

A CSS cursor blink animation provides visual feedback during streaming:

\begin{lstlisting}[language=CSS, caption={Cursor blink keyframe}]
@keyframes blink { 0%, 100% { opacity: 1; } 50% { opacity: 0; } }
\end{lstlisting}

\section{Debug Panel}

When \textit{Debug mode} is enabled (🔍 toggle), each assistant message shows the planner JSON received via the \texttt{planner} SSE event:

\begin{itemize}
  \item Selected tool and confidence
  \item Predicted domain and equipment tags
  \item Retrieval profile
  \item Whether rule-based fallback was used
  \item Reasoning text (\texttt{why\_selected})
\end{itemize}

\section{Key UI Features}

\begin{itemize}
  \item \textbf{LLM status indicator} --- API availability ping on mount; pulsing green dot when connected.
  \item \textbf{Domain filter sidebar} --- restrict retrieval to thermal, electrical, or both.
  \item \textbf{Tool mode selector} --- nine modes with descriptions; \texttt{auto} by default.
  \item \textbf{Evidence panel} --- expandable citation cards with snippet preview, page numbers, relevance score, domain colour.
  \item \textbf{Follow-up suggestions} --- 2 contextual follow-up buttons per response.
  \item \textbf{Explanation level} --- Beginner vs Engineer mode.
\end{itemize}

% ============================================================
\chapter{Evaluation Framework}
% ============================================================

\section{Objectives}

The evaluation framework serves three functions:
\begin{enumerate}
  \item \textbf{Regression detection} --- prevent performance degradation from system changes.
  \item \textbf{Weakness identification} --- isolate which query categories or retrieval patterns fail.
  \item \textbf{Fine-tuning gate} --- establish a pre/post baseline before any model adaptation.
\end{enumerate}

Evaluation runs entirely locally: no annotation crowd, no external embedding APIs, and terminates in $\leq 30$ minutes for all 29 test cases.

\section{Retrieval Metrics}

Let $\mathcal{R}_q^K$ denote the top-$K$ retrieved chunks for query $q$, $\mathcal{E}_q$ the expected source hints, and $\mathcal{K}_q$ the expected keywords.

\subsection{Recall@K}

\begin{equation}
  \text{Recall@}K = \frac{|\{h \in \mathcal{E}_q : \exists c \in \mathcal{R}_q^K,\ h \subseteq \text{text}(c)\}|}{|\mathcal{E}_q|}
\end{equation}

\textbf{Target:} $\geq 0.60$

\subsection{Precision@K}

\begin{equation}
  \text{Precision@}K = \frac{|\{c \in \mathcal{R}_q^K : \exists k \in \mathcal{K}_q,\ k \subseteq \text{text}(c)\}|}{K}
\end{equation}

\textbf{Target:} $\geq 0.50$

\subsection{Mean Reciprocal Rank}

\begin{equation}
  \text{MRR} = \frac{1}{|Q|}\sum_{q \in Q} \frac{1}{\text{rank}_{q,1}}
\end{equation}

\textbf{Target:} $\geq 0.65$

\subsection{Keyword Coverage}

\begin{equation}
  \text{KWCov} = \frac{|\{k \in \mathcal{K}_q : k \subseteq \bigcup_{c \in \mathcal{R}_q}\text{text}(c)\}|}{|\mathcal{K}_q|}
\end{equation}

\textbf{Target:} $\geq 0.60$

\section{Router / Planner Metrics}

\subsection{Tool Selection Accuracy}

\begin{equation}
  \text{ToolAcc} = \frac{1}{|Q|}\sum_{q\in Q}\mathbb{1}[\hat{t}_q = t_q^*] \qquad \text{Target: } \geq 0.80
\end{equation}

\subsection{Domain Accuracy}

\begin{equation}
  \text{DomainAcc} = \frac{1}{|Q|}\sum_{q\in Q}\mathbb{1}[\hat{d}_q \in \{d_q^*,\texttt{"both"}\}] \qquad \text{Target: } \geq 0.85
\end{equation}

\section{Answer Quality Metrics}

\subsection{Keyword Overlap}

\begin{equation}
  \text{KWOverlap} = \frac{|\{k \in \mathcal{K}_q : k \subseteq \text{answer}_q\}|}{|\mathcal{K}_q|} \qquad \text{Target: } \geq 0.50
\end{equation}

\subsection{Completeness Score}

\begin{equation}
  \text{Completeness} = \frac{|\mathcal{T}(\text{reference}_q)\cap\mathcal{T}(\text{answer}_q)|}{|\mathcal{T}(\text{reference}_q)|} \qquad \text{Target: } \geq 0.45
\end{equation}

where $\mathcal{T}(\cdot)$ denotes the set of unique lowercase tokens.

\subsection{Structure Adherence}

Heuristic score $\in \{0, 0.5, 1.0\}$ checking whether the answer conforms to the expected output format (numbered list, diagnostic structure, comparison language, etc.).

\section{Citation Metrics}

\subsection{Citation Field Completeness}

\begin{equation}
  \text{FieldComp} = \frac{1}{|C_q|}\sum_{c\in C_q}\frac{|\text{fields present in }c|}{4} \qquad \text{Target: } \geq 0.90
\end{equation}

Required fields: \texttt{chunk\_id}, \texttt{book\_name}, \texttt{page\_start}, \texttt{relevance\_score}.

\subsection{Source Hint Match}

\begin{equation}
  \text{SrcMatch} = \frac{|\{h\in\mathcal{E}_q : h \subseteq \text{concat}(C_q)\}|}{|\mathcal{E}_q|} \qquad \text{Target: } \geq 0.50
\end{equation}

\section{Composite Score}

\begin{align}
  \text{composite} &= 0.30\,S_r + 0.25\,S_\rho + 0.30\,S_a + 0.15\,S_c \\[4pt]
  S_r &= 0.4\cdot\text{Recall@K} + 0.3\cdot\text{Prec@K} + 0.2\cdot\text{MRR} + 0.1\cdot\text{KWCov} \\
  S_\rho &= 0.6\cdot\text{ToolAcc} + 0.4\cdot\text{DomainAcc} \\
  S_a &= 0.4\cdot\text{KWOverlap} + 0.4\cdot\text{Completeness} + 0.2\cdot\text{Structure} \\
  S_c &= 0.4\cdot\text{CitPresence} + 0.3\cdot\text{FieldComp} + 0.3\cdot\text{SrcMatch}
\end{align}

\begin{table}[H]
\centering
\caption{Composite Score Performance Thresholds}
\begin{tabular}{@{}cll@{}}
\toprule
\textbf{Score} & \textbf{Rating} & \textbf{Action} \\
\midrule
$\geq 0.75$ & Excellent & Production ready \\
$0.60$--$0.75$ & Good & Minor prompt/retrieval tuning \\
$0.45$--$0.60$ & Fair & Systematic issue; investigate retrieval \\
$< 0.45$ & Poor & Re-run ingestion; check Mistral API rate limits \\
\bottomrule
\end{tabular}
\end{table}

\section{Synthetic Dataset}

The 29-case synthetic dataset spans all 8 tool categories:

\begin{table}[H]
\centering
\caption{Synthetic Dataset Composition}
\begin{tabular}{@{}lccc@{}}
\toprule
\textbf{Category} & \textbf{Cases} & \textbf{Thermal} & \textbf{Electrical} \\
\midrule
QA & 8 & 4 & 4 \\
Troubleshoot & 5 & 3 & 2 \\
Opportunity & 4 & 2 & 2 \\
Checklist & 3 & 1 & 2 \\
Explainer & 3 & 2 & 1 \\
Comparison & 2 & 1 & 1 \\
Navigation & 2 & 1 & 1 \\
Summarize & 2 & 1 & 1 \\
\midrule
\textbf{Total} & \textbf{29} & \textbf{15} & \textbf{14} \\
\bottomrule
\end{tabular}
\end{table}

Each case has: \texttt{id}, \texttt{category}, \texttt{query}, \texttt{expected\_tool}, \texttt{expected\_domain}, \texttt{expected\_keywords} (avg.\ 6.3 per case), \texttt{expected\_source\_hints}, \texttt{reference\_answer}, \texttt{expected\_structure}.

\section{Evaluation Script}

\texttt{evaluation/evaluate.py} provides a fully automated CLI pipeline:

\begin{lstlisting}[language=bash, caption={Evaluation CLI}]
python evaluation/evaluate.py \
  --base-url  http://localhost:8000 \
  --dataset   evaluation/synthetic_dataset.json \
  --k         5 \
  --outdir    evaluation/results \
  --run-name  run_001 \
  [--category troubleshoot] \
  [--quiet]
\end{lstlisting}

\noindent Outputs per run:
\begin{itemize}
  \item \texttt{results.jsonl} --- per-case metrics
  \item \texttt{results.csv} --- spreadsheet-ready
  \item \texttt{report.md} --- human-readable with pass/fail table
  \item \texttt{aggregate.json} --- machine-readable global + per-category metrics
\end{itemize}

\section{Latency Targets}

\begin{table}[H]
\centering
\caption{System Latency Budget}
\begin{tabular}{@{}llr@{}}
\toprule
\textbf{Stage} & \textbf{Component} & \textbf{Target (mean)} \\
\midrule
Planner & Mistral API JSON mode & $\leq 3{,}000$ ms \\
Retrieval & Dense + sparse + merge & $\leq 500$ ms \\
Reranking & Cross-encoder & $\leq 600$ ms \\
Generation & Mistral API streaming (first token) & $\leq 5{,}000$ ms \\
\midrule
Total (p50) & End-to-end median & $\leq 10{,}000$ ms \\
Total (p90) & 90th percentile & $\leq 15{,}000$ ms \\
\bottomrule
\end{tabular}
\end{table}

% ============================================================
\chapter{Running and Testing the System}
% ============================================================

\section{Prerequisites}

The following must be installed and running before starting:

\begin{enumerate}
  \item \textbf{Tesseract OCR}: \texttt{sudo apt-get install tesseract-ocr tesseract-ocr-eng}
  \item \textbf{Mistral API Key}: Set \texttt{MISTRAL\_API\_KEY} in \texttt{.env}
  \item \textbf{Python 3.12} with \texttt{.venv} activated and \texttt{backend/requirements.txt} installed
  \item \textbf{Node.js 20} with \texttt{frontend/} dependencies installed
  \item \textbf{PDFs} in the project root
\end{enumerate}

\section{First Run (One-Time Ingestion)}

\begin{lstlisting}[language=bash]
# Quick test (~5-10 min)
./run_ingest.sh --test-run

# Full ingestion (~75-110 min first time; subsequent <20 min)
./run_ingest.sh
\end{lstlisting}

\section{Starting the System}

\begin{lstlisting}[language=bash]
# Terminal 1: Backend
./run_backend.sh          # http://localhost:8000

# Terminal 2: Frontend
cd frontend && npm run dev # http://localhost:3000
\end{lstlisting}

\section{API Smoke Tests}

\begin{lstlisting}[language=bash]
# Health
curl http://localhost:8000/api/health | python3 -m json.tool

# Non-streaming query
curl -X POST http://localhost:8000/api/query \
  -H "Content-Type: application/json" \
  -d '{"query":"What is optimum excess air for a coal boiler?",
       "tool_mode":"auto","explanation_level":"engineer"}'

# SSE streaming
curl -X POST http://localhost:8000/api/stream \
  -H "Content-Type: application/json" \
  -H "Accept: text/event-stream" \
  -d '{"query":"Why is my boiler below 75% efficiency?",
       "tool_mode":"auto","explanation_level":"engineer"}'

# Classify only
curl -X POST http://localhost:8000/api/classify \
  -H "Content-Type: application/json" \
  -d '{"query":"Compare fire tube vs water tube boilers"}'
\end{lstlisting}

\section{Unit Tests}

\begin{lstlisting}[language=bash]
source .venv/bin/activate
# LangGraph tests (18 cases)
PYTHONPATH=backend python -m pytest backend/tests/test_graph.py -v
# OCR cache tests (9 cases)
PYTHONPATH=backend python -m pytest backend/tests/test_ocr.py -v
# All tests
PYTHONPATH=backend python -m pytest backend/tests/ -v
\end{lstlisting}

\section{Automated Evaluation}

\begin{lstlisting}[language=bash]
# Run 29-case evaluation
python evaluation/evaluate.py --run-name run_001

# View results
cat evaluation/results/run_001/report.md
\end{lstlisting}

% ============================================================
\chapter{Future Scope: Fine-Tuning Strategy}
% ============================================================

\begin{keybox}[Scope Constraint]
This chapter presents a design plan only. No training code is implemented. Fine-tuning is positioned as a last-resort optimisation, applied strictly after retrieval, routing, and prompt improvements have been fully exhausted.
\end{keybox}

\section{Key Principle}

\begin{mdframed}
\textbf{RAG handles knowledge. Fine-tuning handles behaviour.}

The BEE manuals are the exclusive source of truth. Fine-tuning shall not be used to inject domain facts into model weights --- that would bypass the RAG grounding and risk hallucination of stale information.
\end{mdframed}

Fine-tuning targets only:
\begin{enumerate}
  \item Tool routing reliability (structured JSON output consistency)
  \item Grounded response style (evidence-backed, never extrapolate)
  \item Output format adherence (checklists, diagnoses, comparisons)
  \item Answer conciseness for engineering-grade queries
\end{enumerate}

\section{Improvement Priority Chain}

Fine-tuning is Step 5 in the improvement chain:

\begin{equation*}
  \underbrace{1.\ \text{Better OCR}}_{\text{ingestion}} \;\longrightarrow\;
  \underbrace{2.\ \text{Better retrieval}}_{\text{chunking, K, RRF}} \;\longrightarrow\;
  \underbrace{3.\ \text{Better routing}}_{\text{planner prompt}} \;\longrightarrow\;
  \underbrace{4.\ \text{Better prompts}}_{\text{system prompts}} \;\longrightarrow\;
  \underbrace{5.\ \text{Fine-tune}}_{\text{last resort}}
\end{equation*}

Proceed to fine-tuning only if composite score stagnates below 0.60 after exhausting all earlier steps.

\section{Training Data Collection}

\subsection{Dataset A --- Router Dataset}

\textbf{Input:} user query \\
\textbf{Output:} complete planner JSON object \\
\textbf{Size target:} 500--1000 examples \\
\textbf{Sources:} expanded synthetic dataset (varied phrasings), production logs with human verification, adversarial negatives (routing failures)

\subsection{Dataset B --- Answer Formatting Dataset}

\textbf{Input:} query + retrieved context chunks + tool mode \\
\textbf{Output:} ideal structured answer \\
\textbf{Size target:} 300--600 examples \\
\textbf{Sources:} reference answers from synthetic dataset, human-corrected production answers

\subsection{Dataset C --- Failure Correction Dataset}

\textbf{Collection trigger:} category composite score $< 0.40$ \\
\textbf{Format:} DPO (Direct Preference Optimisation) preference pairs \\
\begin{equation}
  \text{pair} = (\text{prompt},\ \underbrace{\text{correct grounded answer}}_{\text{chosen}},\ \underbrace{\text{hallucinated/wrong answer}}_{\text{rejected}})
\end{equation}

\section{Fine-Tuning Method}

Full fine-tuning of a 7B model requires $\sim$80 GB VRAM (infeasible on RTX 4050). The only viable approach is QLoRA:

\begin{table}[H]
\centering
\caption{Recommended QLoRA Configuration}
\begin{tabular}{@{}lll@{}}
\toprule
\textbf{Parameter} & \textbf{Value} & \textbf{Rationale} \\
\midrule
Method & QLoRA & 4-bit; fits in 4--6 GB VRAM \\
Rank $r$ & 16--32 & Capacity vs. VRAM balance \\
$\alpha$ & $2r$ & Standard LoRA heuristic \\
Target modules & q\_proj, v\_proj & Attention layers (most impactful) \\
Batch size & 2 & VRAM constraint \\
Gradient accumulation & 16 & Effective batch = 32 \\
Learning rate & 2e-4 & LoRA standard \\
Epochs & 3 & With early stopping \\
Max sequence length & 2048 & Context budget \\
\bottomrule
\end{tabular}
\end{table}

\begin{notebox}[Recommended Tool]
Use \textbf{Unsloth} (\url{https://github.com/unslothai/unsloth}) for QLoRA on consumer GPUs. It provides 2$\times$ faster training and 50\% lower VRAM than vanilla HuggingFace PEFT. A Phase 1 router adapter trains in $\approx$20 minutes on a free Google Colab T4 GPU.
\end{notebox}

\section{Training Phases}

\begin{description}
  \item[Phase 1 --- Router Adapter] Train on Dataset A. Goal: $\text{ToolAcc} \geq 0.92$. Deploy if $\Delta\text{ToolAcc} \geq +0.05$ with no regression.
  \item[Phase 2 --- Answer Formatting Adapter] Train on Dataset B. Goal: structure adherence $\geq 0.90$. Evaluate keyword overlap and completeness.
  \item[Phase 3 --- Failure Correction (optional)] DPO on Dataset C, targeting specific failing categories identified by the evaluation framework.
\end{description}

\section{Deployment Pipeline}

After adapter training and merging:

\begin{lstlisting}[language=bash, caption={Deploy fine-tuned model payload}]
# 1. Convert to GGUF (Q4_K_M quantisation)
python llama.cpp/convert_hf_to_gguf.py ./merged_model \
  --outfile copilot_v2.Q4_K_M.gguf --outtype q4_k_m

# 2. Create Modelfile
cat > Modelfile << 'EOF'
FROM ./copilot_v2.Q4_K_M.gguf
SYSTEM "You are an industrial energy efficiency expert. Answer only from retrieved context."
PARAMETER temperature 0.15
PARAMETER num_ctx 4096
EOF

# 3. Register and deploy securely
# Integrate adapter with Mistral API endpoints or local vLLM

# 4. Update .env
echo "MISTRAL_MODEL=copilot-v2" >> .env
\end{lstlisting}

\section{Evaluation Gate}

Before deploying a fine-tuned model:

\begin{enumerate}
  \item Re-run \texttt{python evaluation/evaluate.py --run-name post\_finetune}
  \item Compare \texttt{aggregate.json} with baseline
  \item Deploy \textbf{only if}:
  \begin{itemize}
    \item $\Delta\text{ToolAcc} \geq 0$ (no regression)
    \item $\Delta\text{composite} \geq +0.02$
    \item $\Delta\text{latency\_p90} \leq +2{,}000$ ms
    \item No category degrades by $> 0.05$
  \end{itemize}
  \item Roll back to base \texttt{mistral-small-latest} if any condition fails.
\end{enumerate}

% ============================================================
\chapter{Configuration and Deployment}
% ============================================================

\section{Environment Variables}

All settings in \texttt{.env}:

\begin{lstlisting}[caption={.env key settings}]
# LLM
LLM_PROVIDER=mistral
MISTRAL_MODEL=mistral-small-latest
MISTRAL_API_KEY=your_api_key_here

# Retrieval
DENSE_TOP_K=20
SPARSE_TOP_K=20
RERANK_TOP_K=8
MAX_CONTEXT_CHUNKS=6
HYBRID_ALPHA=0.6          # 1.0=dense only, 0.0=sparse only
ENABLE_RERANKING=true

# OCR
OCR_DPI=200
FORCE_OCR=false

# Cache
ENABLE_QUERY_CACHE=true
\end{lstlisting}

\section{Docker}

\begin{lstlisting}[language=bash]
docker-compose up --build
\end{lstlisting}

The \texttt{docker-compose.yml} starts the backend service. The frontend can be added as a second service or served separately.

\section{Tech Stack Summary}

\begin{table}[H]
\centering
\caption{Complete Tech Stack}
\begin{tabular}{@{}lll@{}}
\toprule
\textbf{Layer} & \textbf{Technology} & \textbf{Version} \\
\midrule
Orchestration & LangGraph & $\geq$0.2 \\
LLM & Mistral API & cloud \\
OCR & Tesseract & $\geq$5.0 \\
Embeddings & all-MiniLM-L6-v2 (CPU) & local \\
Vector Store & ChromaDB HNSW & 0.5.23 \\
Sparse Search & rank-bm25 (Okapi) & 0.2.2 \\
Reranker & ms-marco-MiniLM-L-6-v2 & local \\
Backend & FastAPI + uvicorn & 0.115 \\
Streaming & sse-starlette & $\geq$2.1 \\
Frontend & Next.js 14 App Router & 14.x \\
Response Cache & diskcache & $\geq$5.6 \\
Testing & pytest + pytest-asyncio & 8.x \\
\bottomrule
\end{tabular}
\end{table}

% ============================================================
\appendix
\chapter{API Request / Response Reference}
% ============================================================

\section{POST /api/query}

\begin{lstlisting}[caption={Request body}]
{
  "query":             "string",
  "tool_mode":         "auto | qa | explainer | ...",
  "domain_filter":     "thermal | electrical | null",
  "explanation_level": "beginner | engineer"
}
\end{lstlisting}

\begin{lstlisting}[caption={Response (abbreviated)}]
{
  "answer":       "string",
  "citations":    [{"chunk_id":...,"book_name":...,"page_start":42,...}],
  "classification": {
    "tool_mode":     "troubleshoot",
    "utility_domain":"thermal",
    "confidence":    0.88,
    "planner_error": false
  },
  "node_latency": {
    "planner_router":   1840,
    "retrieve_dense":     95,
    "retrieve_sparse":    22,
    "merge_results":       8,
    "rerank":            340,
    "answer_generation":5200
  },
  "latency_ms":          7580,
  "follow_up_suggestions":["..."]
}
\end{lstlisting}

\chapter{Test Case Dataset Schema}

\begin{lstlisting}[caption={Synthetic dataset schema (JSON)}]
{
  "id":                    "qa_thermal_001",
  "category":              "qa",
  "query":                 "...",
  "expected_tool":         "qa",
  "expected_domain":       "thermal",
  "expected_keywords":     ["excess air", "boiler", ...],
  "expected_source_hints": ["boiler", "combustion"],
  "reference_answer":      "...",
  "expected_structure":    "direct"
}
\end{lstlisting}

Valid \texttt{expected\_structure} values:
\texttt{direct}, \texttt{structured}, \texttt{structured\_diagnosis}, \texttt{bullet\_list}, \texttt{comparison\_table}, \texttt{checklist}.

\chapter{Evaluation Diagnosis Decision Tree}

\begin{enumerate}
  \item \textbf{Composite $< 0.45$}:
    \begin{enumerate}
      \item Is index empty? $\rightarrow$ run \texttt{./run\_ingest.sh --force-rebuild}
      \item Is Mistral API returning errors? $\rightarrow$ check \texttt{MISTRAL\_API\_KEY} and network connection
    \end{enumerate}
  \item \textbf{Tool accuracy $< 0.70$}:
    \begin{enumerate}
      \item Review planner prompt in \texttt{graph/planner.py}
      \item Check \texttt{planner\_fallback\_rate} --- if high, Mistral API is timing out
      \item Add keywords to rule-based classifier
    \end{enumerate}
  \item \textbf{Recall@K $< 0.40$}:
    \begin{enumerate}
      \item Re-run ingestion with more pages
      \item Check OCR quality: \texttt{python scripts/ocr\_extract.py --quality-report}
      \item Reduce \texttt{HYBRID\_ALPHA} toward 0.5 (more BM25 weight)
      \item Increase \texttt{MAX\_CONTEXT\_CHUNKS} (6 $\rightarrow$ 8)
    \end{enumerate}
  \item \textbf{Keyword overlap $< 0.35$}:
    \begin{enumerate}
      \item Review system prompts in \texttt{generation/prompts.py}
      \item Adjust inference parameters or evaluate larger Mistral models if issues persist
    \end{enumerate}
  \item \textbf{Total p90 $> 20{,}000$ ms}:
    \begin{enumerate}
      \item Verify GPU usage: \texttt{nvidia-smi}
      \item Set \texttt{ENABLE\_QUERY\_CACHE=true}
      \item Reduce \texttt{MAX\_GENERATION\_TOKENS} (1500 $\rightarrow$ 800)
    \end{enumerate}
\end{enumerate}

\end{document}
